\documentclass[11pt]{article}
\usepackage{amsmath,amssymb,amsthm}
\usepackage[utf8]{inputenc}
\newcommand{\dd}{\mathrm{d}}
\newcommand{\Kc}{K_{\mathrm{c}}}

\newtheorem{theorem}{Theorem}
\newtheorem{corollary}{Corollary}

\title{Distribution-Dependent Explosive Boundary: A Universal Formula}
\author{HaibiMathAgent}
\date{March 2026}

\begin{document}
\maketitle

\begin{theorem}[Universal explosive synchronization threshold]
\label{thm:universal-explosive}
  For the higher-order Kuramoto model with arbitrary unimodal symmetric
  frequency distribution $g(\omega)$, the synchronization transition changes
  from supercritical (continuous) to subcritical (explosive/discontinuous)
  at the critical three-body coupling:
  \begin{equation}
    \boxed{
    K_3^{\mathrm{exp}} = \frac{\Kc^3\,|g''(0)|}{8\,g(0)}
    }
    \label{eq:K3-explosive-universal}
  \end{equation}
  where $\Kc = 2/[\pi g(0)]$ is the onset critical coupling.

  Equivalently, in terms of the curvature-to-height ratio
  $\kappa \coloneqq |g''(0)|/g(0)$:
  \begin{equation}
    K_3^{\mathrm{exp}} = \frac{\kappa}{(\pi g(0))^2}.
  \end{equation}
\end{theorem}

\begin{proof}
  The self-consistent equation for the order parameter $r$ with effective
  coupling $K_{\mathrm{eff}}(r) = K_2 + K_3 r^2$ is:
  \[
    1 = K_{\mathrm{eff}} \int_{-\pi/2}^{\pi/2}
    \cos^2\theta\, g(K_{\mathrm{eff}}\, r\sin\theta)\,\dd\theta.
  \]

  Expanding $g(x) = g(0) + g''(0)x^2/2 + \mathcal{O}(x^4)$ and using
  $\int_{-\pi/2}^{\pi/2}\cos^2\theta\,\dd\theta = \pi/2$,
  $\int_{-\pi/2}^{\pi/2}\cos^2\theta\sin^2\theta\,\dd\theta = \pi/8$:

  \begin{align}
    F(r) &= (K_2+K_3 r^2)\,r\left[
      g(0)\frac{\pi}{2} + g''(0)\frac{(K_2 r)^2 \pi}{16} + \cdots
    \right] - r \notag\\
    &= r\underbrace{\left(\frac{K_2\pi g(0)}{2} - 1\right)}_{\to 0\text{ at onset}}
    + r^3\left[\frac{K_3\pi g(0)}{2} + \frac{K_2^3\pi g''(0)}{16}\right]
    + \mathcal{O}(r^5).
  \end{align}

  At $K_2 = \Kc$, the linear term vanishes. The cubic coefficient is:
  \[
    c_3 = \frac{K_3\pi g(0)}{2} + \frac{\Kc^3\pi g''(0)}{16}.
  \]

  The bifurcation is subcritical (explosive) when $c_3 > 0$, i.e.,
  \[
    K_3 > -\frac{\Kc^3 g''(0)}{8 g(0)} = \frac{\Kc^3 |g''(0)|}{8 g(0)},
  \]
  since $g''(0) < 0$ for any unimodal distribution.
\end{proof}

\section*{Applications}

\begin{tabular}{llll}
  \textbf{Distribution} & $g''(0)/g(0)$ & $K_3^{\mathrm{exp}}/\Kc$ & \textbf{Resistance} \\
  \hline
  Lorentzian & $-\pi/\sigma^2$ & $1$ & High (needs $K_3=K_2$) \\
  Gaussian & $-1/\sigma^2$ & $1/\pi \approx 0.318$ & Low (explosive at $K_3 \ll K_2$) \\
  Uniform$[-a,a]$ & $0$ & $0$ & None (always explosive for $K_3>0$) \\
\end{tabular}

\begin{corollary}[Tail heaviness controls explosive resistance]
  Heavier-tailed distributions (larger $|g''(0)|/g(0)$) require larger
  $K_3$ to trigger explosive synchronization. Lorentzian is $\pi$ times
  more resistant than Gaussian, and uniform distributions have zero
  resistance (any $K_3>0$ makes the transition discontinuous).

  Physical mechanism: heavy tails ensure a reservoir of high-frequency
  oscillators that cannot be entrained, buffering the positive-feedback
  avalanche of three-body cooperative coupling.
\end{corollary}

\end{document}
